\textbf{COLUNI 2019 -- Questão 21}\par\medskip \begin{center}\textbf{Brazil museum fire: Funding cuts blamed as icon is destroyed}\end{center}
\begin{flushright}3 September 2018\end{flushright}
\begin{center}(Fonte: <https://ichef.bbci.co.uk/news/976/cpsprodpb/6227/production/_103272152_gettyimages-1026327264.jpg>. Acesso em: 17 set. 2018.)\end{center}

\textbf{Officials in Brazil have blamed lack of funding for a huge fire that has destroyed the country's National Museum.}

One of the largest anthropology and natural history collections in the Americas was almost totally destroyed in Sunday's fire in Rio de Janeiro. This included the 12,000-year-old remains of the skeleton of a woman known as "Luzia" -- the oldest discovered in Latin America. There had also been a string of complaints about the dilapidated state of the museum. "We never had adequate support," its deputy director said after the fire. Experts had warned for years of a serious fire risk to the building.

The fire started on Sunday evening, after the building -- a 19th Century former royal palace -- closed for the day. The cause is not known, but Culture Minister Sergio Sa Leitao was quoted in Brazilian media as saying it may have been ignited by a small paper hot air balloon landing on the roof. By Monday morning the fire was under control and some of the museum's pieces had been rescued. No injuries have been reported but most of the 20 million items the museum contained went up in flames.

The museum contained much history that is now forever lost. Its 20 million artefacts included fossils and dinosaur bones. The jewel in the crown for many visitors was "Luzia". "Luzia is a priceless loss for everyone interested in civilization," museum director Paulo Knauss told the AFP news agency. Using her skull, experts had produced a digital image of her face.

Another popular exhibit was the Bendegó meteorite, weighing more than five tonnes and discovered in Bahía state in the 18th Century. Images taken in the aftermath suggest it is still intact. The building was also home to items covering the centuries from the arrival of the Portuguese in the 1500s to the declaration of a republic in 1889. The ethnology collection had unique pieces from the pre-Columbian era and artifacts from indigenous cultures. Pieces from Greco-Roman times and Egypt were also on display at the museum.

The museum was established in 1818, with the aim of promoting scientific research by making its collection available to specialists. Marcelo Moreira -- a journalist with TV Globo in Rio -- told the BBC: "It's very sad... This museum has a great history. It represents a lot for Brazilian history and Brazilian culture. It's really a big loss for Brazil." In an interview with Globo, the museum's director said it was a "cultural tragedy".

\begin{center}(Disponível em: <https://www.bbc.co.uk/news/world-latin-america-45392668>. Acesso em: 17 set. 2018. Adaptado.)\end{center}

Assinale a afirmativa INCORRETA de acordo com as informações apresentadas no texto:\begin{enumerate}
\item Paulo Knauss disse que o incêndio do Museu Nacional foi uma “tragédia cultural”.
\item O terrível fogo que devastou o museu foi controlado na tarde de terça-feira.
\item Considerando a data de seu estabelecimento, o museu tinha 200 anos.
\item O Ministro da Cultura Sérgio Sá Leitão disse que o fogo pode ter sido iniciado por um balão.\end{enumerate}
